% !TEX program = xelatex
% 元素反应机制详解（7.0 版）—— 共用导言区
% 样式沿用原 Detailed.Explanation.of.Elemental.Reactions.tex（ctexart + 同名宏），并新增元素配色宏与图标宏。
\documentclass[UTF8]{ctexart}
\usepackage{geometry}
\geometry{a4paper,left=1.5cm,right=1.5cm,top=1.5cm,bottom=1.5cm}
\usepackage{array}
\usepackage{booktabs}
\usepackage{longtable}
\usepackage{tabularx}
\usepackage{multirow}
\usepackage[table]{xcolor}
\usepackage{fixcmex}
\usepackage{amsmath}
\usepackage{enumitem}
\usepackage{makecell}
\usepackage{caption}
\usepackage{float}
\usepackage{needspace}
\usepackage{graphicx}
\usepackage[normalem]{ulem}
\usepackage[most]{tcolorbox}
\tcbuselibrary{skins,breakable}
\usepackage{tikz}
\usetikzlibrary{arrows.meta,positioning,shapes.geometric,calc,fit,backgrounds,decorations.pathreplacing}

% ---------- 颜色：七元素 + 月 / 星 ----------
\definecolor{elhydro}{HTML}{3E9BD0}     % 水
\definecolor{elpyro}{HTML}{D8622B}      % 火
\definecolor{elcryo}{HTML}{4FA8C8}      % 冰
\definecolor{elelectro}{HTML}{9A5FC9}   % 雷
\definecolor{elanemo}{HTML}{3FA98A}     % 风
\definecolor{elgeo}{HTML}{C08A18}       % 岩
\definecolor{eldendro}{HTML}{6FA81B}    % 草
\definecolor{ellunar}{HTML}{5A6FD8}     % 月（月曜体系）
\definecolor{elstellar}{HTML}{C2569E}   % 星（星烁体系）
\definecolor{srcgray}{HTML}{6E6E6E}     % 引注灰
\definecolor{termcolor}{HTML}{1F6F6B}   % 概念名词
\definecolor{rxamp}{HTML}{C2611E}       % 增幅类
\definecolor{rxtrans}{HTML}{1E6FA8}     % 剧变类
\definecolor{rxcryst}{HTML}{8A6A1E}     % 结晶类
\definecolor{rxcata}{HTML}{2F8B4C}      % 激化类
\definecolor{rxlunar}{HTML}{4A5CC4}     % 月曜类
\definecolor{rxstellar}{HTML}{A8559B}   % 星烁类
% 反应名配色：取 bwiki《模板:颜色》各反应渐变描边/元素色（LaTeX 不做渐变，取可读的元素身份色）
\definecolor{rxlc}{HTML}{9245E6}       % 月感电（bwiki 雷元素色）
\definecolor{rxlb}{HTML}{619E7F}       % 月绽放（bwiki 草元素色）
\definecolor{rxlcrys}{HTML}{A99847}    % 月结晶（bwiki 岩元素色）
\definecolor{rxssc}{HTML}{8784E0}      % 星超导（bwiki 独特蓝紫，非冰）
\definecolor{rxssw}{HTML}{2BC5A2}      % 星扩散（bwiki 风元素色）
\definecolor{bwikiDesc}{HTML}{BA7920}   % bwiki 描述色（月兆/辉映/伪秘之帷/擢升等）
\definecolor{prism}{HTML}{6D8DFB}        % 星辉棱晶（游戏内偏白蓝色）

% ---------- 元素名 / 图标 ----------
\newcommand{\elcol}[1]{\csname elc@#1\endcsname}
\expandafter\def\csname elc@水\endcsname{elhydro}
\expandafter\def\csname elc@火\endcsname{elpyro}
\expandafter\def\csname elc@冰\endcsname{elcryo}
\expandafter\def\csname elc@雷\endcsname{elelectro}
\expandafter\def\csname elc@风\endcsname{elanemo}
\expandafter\def\csname elc@岩\endcsname{elgeo}
\expandafter\def\csname elc@草\endcsname{eldendro}
\expandafter\def\csname elc@月\endcsname{ellunar}
\expandafter\def\csname elc@星\endcsname{elstellar}
\expandafter\def\csname elc@物理\endcsname{srcgray}
\expandafter\def\csname elc@岩元素\endcsname{elgeo}
% 容错：未登记的元素键直接排原文/不插图，避免 \csname 报错
\newcommand{\elname}[1]{\ifcsname elc@#1\endcsname\textcolor{\elcol{#1}}{#1}\else #1\fi}
\newcommand{\element}[1]{\ifcsname elc@#1\endcsname\textcolor{\elcol{#1}}{[#1]}\else [#1]\fi}

\newcommand{\iconfile}[1]{\csname iconf@#1\endcsname}
\expandafter\def\csname iconf@水\endcsname{icons/water.png}
\expandafter\def\csname iconf@火\endcsname{icons/fire.png}
\expandafter\def\csname iconf@冰\endcsname{icons/ice.png}
\expandafter\def\csname iconf@雷\endcsname{icons/electro.png}
\expandafter\def\csname iconf@风\endcsname{icons/anemo.png}
\expandafter\def\csname iconf@岩\endcsname{icons/geo.png}
\expandafter\def\csname iconf@草\endcsname{icons/dendro.png}
\newcommand{\icon}[1]{\ifcsname iconf@#1\endcsname\raisebox{-0.22ex}{\includegraphics[height=1.05em]{\iconfile{#1}}}\else\relax\fi}

% ---------- 反应名 / 概念名词 / 引注 / 存疑 ----------
\newcommand{\rxAmp}[1]{\textcolor{rxamp}{\textbf{#1}}}
\newcommand{\rxTrans}[1]{\textcolor{rxtrans}{\textbf{#1}}}
\newcommand{\rxCryst}[1]{\textcolor{rxcryst}{\textbf{#1}}}
\newcommand{\rxCata}[1]{\textcolor{rxcata}{\textbf{#1}}}
\newcommand{\rxLunar}[1]{\textcolor{rxlunar}{\textbf{#1}}}
\newcommand{\rxStellar}[1]{\textcolor{rxstellar}{\textbf{#1}}}
% 反应名着色宏：元素色填充 + 同色浅投影（模拟游戏内反应名高光）；体系名（月曜／星烁）用 \rxLunar／\rxStellar
\newcommand{\rxLC}[1]{\leavevmode\rlap{\kern0.4pt\lower0.4pt\hbox{\textcolor{rxlc!30}{\textbf{#1}}}}\textcolor{rxlc}{\textbf{#1}}}
\newcommand{\rxLB}[1]{\leavevmode\rlap{\kern0.4pt\lower0.4pt\hbox{\textcolor{rxlb!30}{\textbf{#1}}}}\textcolor{rxlb}{\textbf{#1}}}
\newcommand{\rxLCrys}[1]{\leavevmode\rlap{\kern0.4pt\lower0.4pt\hbox{\textcolor{rxlcrys!30}{\textbf{#1}}}}\textcolor{rxlcrys}{\textbf{#1}}}
\newcommand{\rxSSC}[1]{\leavevmode\rlap{\kern0.4pt\lower0.4pt\hbox{\textcolor{rxssc!30}{\textbf{#1}}}}\textcolor{rxssc}{\textbf{#1}}}
\newcommand{\rxSSW}[1]{\leavevmode\rlap{\kern0.4pt\lower0.4pt\hbox{\textcolor{rxssw!30}{\textbf{#1}}}}\textcolor{rxssw}{\textbf{#1}}}
% 概念名词着色宏：按 bwiki 配色；bwiki 未着色的术语用 \term（teal）
\newcommand{\termDesc}[1]{\textcolor{bwikiDesc}{\textbf{#1}}}    % bwiki 描述色
\newcommand{\termElectro}[1]{\textcolor{rxlc}{\textbf{#1}}}      % bwiki 雷色（复用 rxlc）
\newcommand{\termDendro}[1]{\textcolor{rxlb}{\textbf{#1}}}      % bwiki 草色（复用 rxlb）
\newcommand{\termAnemo}[1]{\textcolor{rxssw}{\textbf{#1}}}      % bwiki 风色（复用 rxssw）
\newcommand{\termGeo}[1]{\textcolor{rxlcrys}{\textbf{#1}}}     % bwiki 岩色（复用 rxlcrys）
\newcommand{\termPrism}[1]{\textcolor{prism}{\textbf{#1}}}     % 星辉棱晶
\newcommand{\term}[1]{\textcolor{termcolor}{\textbf{#1}}}
\newcommand{\src}[1]{{\footnotesize\textcolor{srcgray}{［#1］}}}
% 交叉引用：章节一律带 § 前缀（用户口径：行文交叉引用要规范），且 hyperref 下可点击
\newcommand{\sref}[1]{\S\ref{#1}}
\newcommand{\cunyi}{\textcolor{red}{\textbf{【存疑】}}}
\newcommand{\thickline}{\Xhline{1.2pt}}

% ---------- 版式微调（沿用原文档观感；标题实现由 titlesec 改为 ctexset） ----------
% 变更原因：titlesec 与 hyperref 组合在分页处锚点不可靠（hyperref 手册 11.1.33 明言不支持），
% 会造成目录点击「跳到前一页」；ctexset 是 ctex 原生实现，锚点与 hyperref 兼容。
% 观感对齐原 titlesec 配置：格式 \normalfont\Large\bfseries、编号后 1em、前后距 1ex/0.5ex；
% afterindent=false 复刻 titlespacing* 的「标题后首段不缩进」；fixskip=true 抑制额外间距。
\ctexset{
  section = {
    format      = \normalfont\Large\bfseries,
    aftername   = \hspace{1em},
    beforeskip  = 2.2ex,
    afterskip   = 1.4ex,
    afterindent = true,
    fixskip     = true,
  },
  subsection = {
    format      = \normalfont\large\bfseries,
    aftername   = \hspace{1em},
    beforeskip  = 1.8ex,
    afterskip   = 1.6ex,
    afterindent = true,
    fixskip     = true,
  },
  subsubsection = {
    format      = \normalfont\small\bfseries,
    aftername   = \hspace{1em},
    beforeskip  = 1.4ex,
    afterskip   = 1.2ex,
    afterindent = true,
    fixskip     = true,
  },
}
\renewcommand{\arraystretch}{0.9}
\setlength{\tabcolsep}{4pt}
% 列表更紧凑（批注：相干内容行距太大）；一级列表整体缩进到正文首行缩进处
% （批注：「这段属于【反应进程】，却没有缩进，反而比标签还往左」）
\setlist{topsep=0.1ex,itemsep=0ex,parsep=0pt,partopsep=0pt}
\setlist[itemize,1]{labelindent=2em,leftmargin=*}
\setlist[enumerate,1]{labelindent=2em,leftmargin=*}
% 图表标题字号小于正文（批注：图例的文字不能和正文一样）
% 批注4：「图例标题、表的标题都要用正楷」→ 标题一律楷体（与来源区同一字体）
\newCJKfontfamily\kaifont{simkai.ttf}
\DeclareCaptionFont{kaicap}{\kaifont\let\textbf\relax}
\captionsetup{font={small,kaicap}}
% 公式块前后留白收紧，避免"相干内容行距太大"
\AtBeginDocument{%
  \setlength{\abovedisplayskip}{0.8ex plus 0.2ex minus 0.2ex}%
  \setlength{\belowdisplayskip}{0.8ex plus 0.2ex minus 0.2ex}%
  \setlength{\abovedisplayshortskip}{0.4ex}%
  \setlength{\belowdisplayshortskip}{0.4ex}%
}

\xeCJKDeclareCharClass{CJK}{"2460 -> "24FF}
% 说明：①②③④ 等带圈数字在 Latin Modern 中无字形，XeLaTeX 会静默丢字；
% 交给中文字体渲染后 Missing character 归零。

% ---------- 来源说明块 / 表格行线（第六轮批注口径）----------
% 来源块：楷体、小于正文、上方一条灰横线（全局统一，不再与正文样式混用）
% （楷体族 \kaifont 已在标题设置处定义，图表标题与本区共用）
\newcommand{\srcrule}{\par\vspace{1.5ex}{\color{black!30}\hrule height 0.4pt}\vspace{0.35ex}}
\newenvironment{srcnote}
  {\par\srcrule\begingroup\kaifont\small\linespread{1.05}\selectfont
   \let\textbf\relax           % 楷体无粗体字形，避免 font-shape 回退告警
   \setlength{\parindent}{0pt}\setlength{\parskip}{0.2ex}}
  {\par\endgroup\vspace{0.3ex}}
% 注与来源说明：按批注「注不应写在正文中，应写在页脚」「来源区都是放到那里的」移到页脚。
% 用法：表格内用 \srcmark（浮动体内不能放脚注正文），紧随 \end{table} 之后用 \srctext{...}。
% \srcmark 用 \DeclareRobustCommand 包装：\footnotemark 是脆弱命令，直接写进 \caption 会报
% 「Argument of \caption@ydblarg has an extra }」（实测）。
\DeclareRobustCommand{\srcmark}{\footnotemark}
\newcommand{\srctext}[1]{\footnotetext{{\kaifont\footnotesize\textcolor{srcgray}{\let\textbf\relax #1}}}}
% 来源说明块：同样落在页脚，但不编号——批注4「这些应放在页末（页脚）……一条灰线就够了」。
% 不编号还避开一个坑：\footnotetext 用的是“当前”脚注计数（它本身不 step counter），
% 同一页若有两个 mark 与两段注文交错，两段注文会印出同一个号；来源块不带号即无此问题。
\newcommand{\srcfoot}[1]{%
  \begingroup
  \renewcommand{\thefootnote}{}%
  \footnotetext{{\kaifont\footnotesize\textcolor{srcgray}{\let\textbf\relax #1}}}%
  \endgroup}
% 表格行线：每行一条浅灰细线（0.3pt），不改变三线表的顶/底粗线
\newcommand{\rowline}{\hline}
% 公式块的小标签（如「增幅反应」）：批注4「放到下一页吧，放这里的话，下一页的内容开头
% 看完之后，要往上翻一页才知道是属于增幅反应的公式」→ 标签与其后公式块不许被分页拆开。
\newcommand{\rxlabel}[1]{\needspace{6\baselineskip}\indent\textbf{#1}}

% ---------- 阅读性组件：结论框 / 公式框 / 存疑框 ----------
\newtcolorbox{keybox}[1]{colback=ellunar!4!white,colframe=ellunar,title={#1},fonttitle=\bfseries,
  boxrule=0.7pt,arc=1.5mm,left=3pt,right=3pt,top=2pt,bottom=2pt,breakable}
\newtcolorbox{rxnewbox}[1]{colback=elstellar!4!white,colframe=elstellar,title={#1},fonttitle=\bfseries,
  boxrule=0.7pt,arc=1.5mm,left=3pt,right=3pt,top=2pt,bottom=2pt,breakable}
\newtcolorbox{warntbox}[1]{colback=red!3!white,colframe=red!55!black,title={#1},fonttitle=\bfseries,
  boxrule=0.7pt,arc=1.5mm,left=3pt,right=3pt,top=2pt,bottom=2pt,breakable}
\newtcolorbox{fmbox}{colback=black!3!white,colframe=black!35,boxrule=0.5pt,arc=1.5mm,
  left=3pt,right=3pt,top=3pt,bottom=3pt,breakable}
\newcommand{\hdr}[1]{\textbf{#1}}
\newcommand{\rowhdr}{\rowcolor{black!8}}

% ---------- 超链接：目录/书签可点击跳转（hidelinks 不改变现有观感） ----------
\usepackage[hidelinks,bookmarksnumbered=true]{hyperref}
\hypersetup{pdftitle={元素反应机制详解（7.0版）}}
